% Transcription — fonds Alexandre Grothendieck, shelfmark 156-7, batch 5 (pages 81-100).
% Established from the facsimile archives/batches/156-7/batch-05.pdf (a copy of
% 156-7.pdf fetched through the project's relay; no cover sheet in the batch,
% so PDF page k = folder page 80+k), per the skill transcribe-grothendieck. The
% folder is 113 pages in six batches, transcribed in parallel; this is the
% fifth. The second draft (folder 156-6, GF VI) and batches 3 and 6 of this
% folder were read for notation; batch 4 (pages 61-80, where the supports,
% the lattice Σ_𝓜, supp / cosupp and the condition « Supp_P 1, 2 » of p. 80
% must be introduced) did not exist when this batch was written, so those
% notions are taken as he writes them.
%
% Pass: Opus 5.5 (claude-opus-5-5), 2026-09-26 — first pass, unchecked against the pages by a human.
%
% Pages skipped: none. Pages summarised: none. All twenty leaves carry his
% mathematics in ink, fast drafting, and all are transcribed. Page 100 is
% written on the back of an unrelated typed leaf; only his notes are given.
%
% His pagination 81-100 stands at the head of each leaf and coincides with the
% archivists' (noted on page 81; the « 84 » is written over another digit).
% Dates in his hand: « 26.6 » (read so; « 26 » underlined), opening the long
% slanted note in the top-left margin of page 86, given as \marginal with a
% \note; « 26 juin », underlined, opening the second paragraph of page 99,
% set as \subsection with a \note. The page-86 margin was written on 26 June
% beside text that precedes the 26 June heading of page 99.
%
% Numbered runs and labels, recorded where they stand: At supp 1 (p. 86) and
% At supp 1 (spéc) (p. 87); a label « (δ) » (read so) on page 85, recalled
% on page 86; Lemme 1 (?), Lemme 2 a) b) (pp. 89-91); Théorème (a)-(d)
% (circled letters, p. 93); Lemme 3 (p. 94), Lemme 4 (p. 95); on pages
% 96-97 an abstract « Lemme » with conditions a), b) and its own Lemmes 1-3,
% then refuted (p. 98); an unnumbered Lemme and Corollaire (pp. 99-100),
% (a), (b) circled on p. 100. Axioms cited: At D, At L 1-3 (L1 « de
% divisibilité »), At D° (?) (p. 89). No formula numbers (1.x) in this
% batch. Internal page references: « cf p. 95 » (margin of p. 91), « cor.
% p. 86 » (p. 99), « cf p. 80, condition Supp_P 1 » (p. 100). No « GF IV /
% GF VI, p. N » citation occurs on these pages.
%
% What the batch is about. Supports and the « interior » X° of a multistrate.
% Pages 81-83 test the formalism of supports on quasi-ensemblist examples:
% quasi-simplexes of a space P with |X| ⊂ P, the « ficelle à deux brins »,
% whose interior is ]0,1[ × ]0,1[, and more generally I × I' for two totally
% ordered sets with extremities: for z = (x, x') the complement B_z has
% |B_z1 ∩ B_z2| ⊊ |B_z1| ∩ |B_z2|, |A| ⊊ |supp A|, and the modular
% distributivity of ∩ over ∨ fails. Conclusion-Scholie (p. 83): the
% formalism of supports is good provided one does not use it to compare
% supports of non-compatible strata. Pages 84-86: for a figure F with
% F̃ = I, the elements of I as « open strata »; the definition
% X° = supp X ∩ cosupp ∂X; the wish supp X = Sup_Σ Y°; a provisional
% heuristic criterion (cancelled); « (δ) » supp omb(X)° ⊂ X° from At L2,
% At L3; the axiom At supp 1 (Y ∈ X° ⇒ Y ≪ X), whence a Corollaire
% X° = {Y ∈ Omb(X) | Y_∂X = ∅}; the 26 June margin finds the axiom
% artificial (not satisfied by the « gros ateliers », nor by the piecewise
% linear ones). Page 87: At supp 1 (spéc). Pages 88-91: setting supports
% aside, Omb°°(X), X° ∥ Y° (cancelled), then the Définition of X° through
% Omb and ∂X (p. 89), Lemme 1 (X ∥ Y iff all X'° ∥ Y'°), Lemme 2 on
% interior refinement ≪̊ (Y ≪̊ X iff Y° ⊂ X°; Y ∈ X° iff Y'° ⊂ X° for all
% strata Y'), and an NB: X ≠ Y ⇒ X° ∩ Y° = ∅. Pages 91-95: for a finite
% closed family Φ = F̃ of pairwise compatible multistrates, S_A = Supp(⋃ X°)
% for A ⊂ F̃, a map 𝔓(F̃) → Σ_𝓜, and a Théorème (injective, an iso onto an
% ordered subset, commuting with Sup and Inf, S_{A∖B} = S_A ∩ ∁S_B, S_A =
% supp G for A closed); proof of (d) and (a) through Lemmes 3 and 4
% (X ≬ Y, X ≠ Y ⇒ X° ∥ Y°). Pages 96-98: the theorem as a case of an abstract
% Lemme on a set with a symmetric antireflexive ∥ and a family (A_i) of
% non-empty pairwise disjoint parts; Lemmes 1-3; then « Mais c'est faux ! »:
% a graph A' — A — B — B' refutes Lemme 3, « Rideau… ». Pages 98-99 hope the
% formula holds under At L1. On 26 June (p. 99) he takes X° = Supp omb(X)°
% instead, and states a Lemme (F ≼ X, F̃° = φ⁻¹({X}): X° = Sup X'°) and a
% Corollaire (p. 100), checked with a support function X ↦ |X| valued in P:
% (a) |X°| = |X| ∖ |∂X|.
%
% Reading hazards fixed once here, aligned with 156-6 and batches 3 and 6:
% 𝔉 \mathfrak{F}, 𝓜 \mathcal{M}, ℒ \mathcal{L}; disjointness \parallel,
% its negation with a stroke \nparallel, with a bar over it
% \overline{\parallel}; compatibility \between; interior refinement ≪ with a
% small circle above \overset{\circ}{\ll}, with two circles
% \overset{\circ\circ}{\ll}; subdivision \preccurlyeq; « supp », « Supp »,
% « cosupp », « Omb », « omb » in roman (he writes both supp and Supp and
% both are kept); Σ_𝓜 \Sigma_{\mathcal{M}}; his complement ∁ (a bold C)
% \complement; the induced part « Y_∂X » as a subscript, the restriction
% « Y'|Y_∂X » with a bar, as written. Interior X° is X^{\circ}; Omb° and
% Omb°° as he writes them. The little scheme at the top of page 90 is set in
% tikz-cd with ≤ and ≪ riding headless lines; the drawings of pages 81, 82
% and 98 are described in French notes. His underlining is \emph; circled
% letters are set in parentheses with a note.
%
% Hard places, read only in part: most of page 81 (fast prose on
% quasi-simplexes and ficelles) and its left margin; the end of the
% Conclusion-Scholie (p. 83); the slanted margins of pages 86, 89, 90, 94
% and 100; the second half of page 98. Variances kept and flagged: « Y'_∂X »
% for Y'_∂Y in the proof of Lemme 4 (p. 95); « supp A = cosupp {A', B} = A' »
% and « A ∪ B = ∁A ∩ ∁B = ∅ » (p. 98); « Z' < X ou Z' < Y » (p. 88).

\documentclass[11pt,a4paper]{article}
\input{../preamble/grothendieck}

\folder{156-7}
\batch{5}
\pages{81}{100}
\dating{1986}
\watermark{Édition de démonstration}
\foldertitle{[Chapitre] VII. Analysis situs (troisième mouture) : notes manuscrites (23-26/06/1986)}

\begin{document}

\page{81}
\note{pagination de sa main 81 à 100 en tête des feuillets, qui coïncide avec celle des archivistes}

\emph{NB} Ceci \uncertain{détermine} $P$ et l'application $A \mapsto |A|$ à iso \uncertain{canonique} près. Je \uncertain{prévois} \ill{} que les conditions (i$_{\mathcal{M}}$) (ii$_{\mathcal{M}}$) (iv$_{\mathcal{M}}$), (v$_{\mathcal{M}}$) (cf At L 2) seront satisfaites \uncertain{automatiquement}\ldots
\note{la ligne « Je prévois \ldots\ automatiquement\ldots » est écrite en plus petit entre les lignes, au-dessus et au-dessous de « à iso canonique près »}

Je prends \uncertain{maintenant} \ill{}, en dehors des \ill{} des ateliers quasi-ensemblistes, où on a \ill{} \ill{} \ill{} \ill{} espace topologique $P$, et on définit $\mathcal{L}$ (« lieux ») comme l'ens.\ des parties compactes contractiles disons \add{de $P$}. On construit des « quasi-simplexes » \add{de \ill{}} de dimension \ill{} \uncertain{conditions} \ill{}, \ill{} \ill{} \ill{} \ill{} (\ill{} pour 1) par des conditions ad hoc, \ill{} définissent \ill{} quasi-ensembliste \ill{} ficelles. [Dans le cas \ill{} $\leq 1$, \ill{} \ill{} \ill{} définition \ill{} forme \ill{} \ill{} \ill{}.] On définit $|X|$ de façon \uncertain{évidente}. Tout marche ! Mais on fait attention que si \ill{} \ill{} aux lieux qui \ill{} \ill{} sur une ficelle donnée, et à la fonction \ill{}, la condition Supp$_{P}$ \uncertain{1, 2} vue ci-dessus, n'est plus satisfaite en général.
\note{paragraphe de prose rapide, lu par fragments ; « lieux » est écrit entre guillemets au-dessus de « $\mathcal{L}$ », « de $P$ » au-dessus de « disons », « quasi-simplexes » entre guillemets. La condition « Supp$_P$ » renvoie au lot 4 de ce dossier (cf.\ page 100 : « p.~80, condition Supp$_P$ 1 »)}

\emph{Exemple} Je prends la « ficelle à deux brins »
\drawing{dessin : un fuseau fermé, deux arcs joignant un point marqué $0$ à gauche à un point marqué $1$ à droite ; une croix $x$ sur l'arc supérieur ; à côté, « (\uncertain{saisons} \ill{}) » ; sous le dessin, « (un \ill{} $x$ de $]0,1[$) »}
\ill{} « \ill{} » de l'un des lieux \ill{} et l'autre brin — ils s'identifient donc \ill{} l'un \ill{} produit
\[
  I^{\circ} = \left]0, 1\right[ \times \left]0, 1\right[ .
\]
\marginal{Attention, ici $x$ \ill{} la fonction \ill{} satisfait \ill{} Supp$_{P}$ 1}
\note{note oblique dans la marge gauche, en face de l'exemple, lue en partie ; un trait la relie à la croix $x$}

\page{82}
Plus généralement, on peut prendre deux \uncertain{treillis} \add{totalement ordonnés} $I, I'$, \struck{\ill{}} (\uncertain{totalement} \ill{}) \add{avec extrémités} prendre le produit $I \times I'$, les \uncertain{ayant} un plus petit et un plus grand élément (« origine » et « extrémité »), \uncertain{et} pour \ill{}
\[
  z = (x, x') \in J^{\circ} = I \times I'
\]
on ait que
\[
  A_z = \lbrace z \rbrace
\]
est un ensemble support, \uncertain{l'ens.}\ des \struck{\ill{}} \add{strates} de la figure correspondant\supplied{e}
\struck{\ill{}}
\[
  [z'', 1] \ \text{pour}\ z'' > z , \qquad [0, z'] \ \text{pour}\ z' < z .
\]
\note{« avec extrémités » est écrit au-dessus d'un mot biffé ; « strates » au-dessus d'un mot biffé}

Considérons la fonction \uncertain{support} \add{\uncertain{naturelle}} \struck{\ill{}} à valeurs dans $I \amalg I' \amalg \lbrace 0 \rbrace \amalg \lbrace 1 \rbrace$ : si $z = (x, x')$, $x \in I$,
\[
  |z| = \lbrace x, x' \rbrace , \qquad |0| = \lbrace 0 \rbrace , \qquad |1| = \lbrace 1 \rbrace .
\]
\note{la réunion disjointe est écrite au-dessus de biffures ; lecture douteuse}
Le complément de $A_z$ est
\[
  B_z = \lbrace z' \in J \mid z' < z \ \text{ou}\ z' > z \rbrace
\]
\struck{On} \uncertain{On} \ill{} \ill{} \ill{} que
\struck{$B_{z} \cap B_{z_2}$}
\[
  |B_z| = J \smallsetminus |z|
\]
\ill{} il se dit
Mais si $z_1, z_2$ ne sont pas \uncertain{comparables}, on a
\[
  |B_{z_1} \cap B_{z_2}| \subsetneqq |B_{z_1}| \cap |B_{z_2}| .
\]
Il manque $]z_1, z_2[$ sur $I$ et $]z'_2, z'_1[$ sur $I'$ \ill{} \uncertain{de dedans} \ill{}
\[
  [x_1, x_2] \cup [x'_2, x'_1] = |A_{z_1} \vee A_{z_2}| \supsetneqq |A_{z_1}| \cup |A_{z_2}| = \lbrace x_1, x_2, x'_1, x'_2 \rbrace
\]
\drawing{dessin à gauche : un fuseau de $0$ à $1$ ; sur l'arc supérieur les points $x_1$, $(x_3)$, $x_2$, sur l'arc inférieur $x'_2$, $(x'_3)$, $x'_1$ ; les segments $[x_1, x_2]$ et $[x'_2, x'_1]$ sont repassés en gras ; à côté, « $z_1 = (x_1, x'_1)$ », « $z_2 = (x_2, x'_2)$ »}

\page{83}
Soit d'autre part
\[
  A = \lbrace z_1, z_2 \rbrace
\]
on aura donc
\[
  |A| = |z_1| \cup |z_2| = \lbrace x_1, x_2, x'_1, x'_2 \rbrace
\]
mais $\mathrm{supp}\, A = A_{z_1} \vee A_{z_2}$ donc
\[
  |\mathrm{supp}\, A| = [x_1, x_2] \cup [x'_1, x'_2]
\]
\struck{ils sont} donc
\[
  |A| \subsetneqq |\mathrm{supp}\, A| .
\]
Considérons $z_3 = (x_3, x'_3)$ avec $x_1 < x_3 < x_2$, $x'_2 < x'_3 < x'_1$, et
\[
  A_{z_3} \cap (A_{z_1} \vee A_{z_2}) = A_{z_3}
\]
\[
  (\underbrace{A_{z_3} \cap A_{z_1}}_{\emptyset}) \vee (\underbrace{A_{z_3} \cap A_{z_2}}_{\emptyset}) = \emptyset
\]
donc on n'a pas la propriété modulaire \ill{} distributivité de $\wedge$ par rapport à $\vee$ !

\emph{Conclusion-Scholie} Le formalisme des supports est bon, pourvu qu'on \struck{\ill{} \ill{}} ne s'en serve pas \ill{} comparer les supports de strates et des figures \ill{} compatibles, ou \struck{\ill{} \ill{}} \add{qui ne sont} \ill{} d'être avec deux multistrates qui ne puissent \ill{} (Tel est le cas \ill{} \ill{} des \ill{} subdivisions \ill{} quasi-ensemblistes, en \ill{} \ill{} \ill{} \ill{}, la définition \uncertain{formelle} \ill{} \ill{}.)
\note{le paragraphe est marqué en marge gauche de deux traits verticaux ; ses dernières lignes, serrées au bas de la page, ne sont lues que par fragments}

\page{84}
\note{le « 84 » de sa main est écrit en surcharge sur un autre chiffre}
C'est cette intuition que je voudrais maintenant essayer de cerner de plus près.

Soit $\mathfrak{F}$ un atelier, $F$ une figure, $\widetilde{F} = I$ l'ens.\ ordonné de ses strates. \struck{Dans le cas \ill{} \ill{} $F = \mathfrak{F}$} Intuitivement, on identifie les éléments de $I$ à des « strates ouvertes » et on veut identifier $\mathfrak{P}(I)$ à \struck{quelque} un ensemble de supports (réunions de « strates ouvertes »).
\note{au-dessus de ce paragraphe, un trait isolé sépare la première phrase du reste}

Soit $X$ une strate. On a envie de donner un sens intrinsèque à « $X \smallsetminus \partial X$ ». L'idée est de \struck{prendre} définir
\[
  X^{\circ} = \mathrm{supp}_{\mathcal{M}} X \cap \mathrm{cosupp}_{\mathcal{M}} \partial X \qquad (\subset \mathcal{M})
\]
où je rappelle
\[
  \partial X = \lbrace Y \in \mathcal{M} \mid Y < X \rbrace = \widetilde{X} \smallsetminus \lbrace X \rbrace
\]
On aimerait d'ailleurs vérifier que
\[
  \mathrm{supp}\, X = \mathop{\mathrm{Sup}_{\Sigma}}_{Y \in \widetilde{X}} Y^{\circ} \qquad \text{Sup dans}\ \Sigma_{\mathcal{M}}
\]
i.e., prenant \uncertain{les} supports complémentaires
\[
  \mathrm{cosupp}\, X = \bigcap_{Y \in \widetilde{X}} \complement_{\mathcal{M}}(Y^{\circ})
\]
\note{devant $\complement_{\mathcal{M}}$, un « cosupp » biffé}

\page{85}
i.e.\ $\forall\, Z \in \mathcal{M}$
\[
  Z \parallel X \Longleftrightarrow \forall\, Y \in \widetilde{X},\ Z \parallel Y^{\circ}
\]
Bien sûr, comme $Y^{\circ} \subset \mathrm{supp}\, Y \subset \mathrm{supp}\, X$ et $Z \parallel X \Leftrightarrow Z \parallel \mathrm{supp}\, X$, l'implication $\Longrightarrow$ est claire. \struck{Donc} Il s'agit de voir si l'inverse est vrai. \struck{Il faudrait donc analyser la condition $Z \parallel Y^{\circ}$}

\struck{Je proposerais le critère heuristique provisoire}

\struck{\emph{Critère} $Z \nparallel Y^{\circ} \Longleftrightarrow \exists$ subdivision $Z'$ de $Z$, et une sous-figure $Z'_0$ de $Z'$, telles que $Z'_0 \ll Y$ et « $Z' \smallsetminus Z'_0$ » $\parallel Y$}
\note{le critère est encadré et barré de trois longs traits obliques ; dans « $Z'_0 \ll Y$ » un signe biffé précède $Y$}

Je voudrais comparer $X^{\circ}$ et $\mathrm{supp}\, \mathrm{omb}(X)^{\circ}$. Je dis que
\[
  (\delta) \qquad \mathrm{supp}\, \mathrm{omb}(X)^{\circ} \subset X^{\circ}
\]
ce qui revient (puisqu'on sait \add{en prenant At L 2} que \struck{c'est} \ill{}) contenu dans $\mathrm{supp}(\mathrm{omb}(X)) = \mathrm{supp}\, X$, à voir que $\mathrm{supp}$ est contenu dans $\mathrm{cosupp}\, \partial X$, i.e.\ \struck{omb} $\mathrm{omb}(X)^{\circ} \parallel \partial X$, ce qui \struck{Compatibles At L 2} est At L 2) et At L 3. \struck{Ainsi} A-t-on égalité ?
\note{le label « ($\delta$) » est lu ainsi sans certitude (un $8$ ou un astérisque ne sont pas exclus) ; il est rappelé page 86. « en prenant At L 2 » est écrit au-dessus de la ligne}

\page{86}
\marginal{26.6 Finalement, cet axiome me \uncertain{semblait} un peu artificiel, il n'est pas satisfait par les « gros ateliers » \ill{} \ill{} \ill{} \ill{} \ill{} \ill{} \ill{} \ill{} \ill{} « artificiel » ! Mais il n'est pas satisfait non plus pour l'atelier linéaire par morceaux \ill{} \ill{} topologiques\ldots}
\note{longue note écrite en oblique dans l'angle supérieur gauche, ouverte par la date « 26.6 » (lue ainsi ; « 26 » souligné) ; elle vise l'axiome At supp 1 énoncé en dessous. Lue en partie seulement : elle chevauche l'énoncé et plusieurs mots en sont recouverts}

\ill{} \ill{} d'introduire un axiome convenable :

\emph{At supp 1} Soient $X, Y \in \mathcal{M}$, supposons que
\[
  Y \in X^{\circ} \overset{\mathrm{def}}{=} \mathrm{supp}\, X \cap \mathrm{cosupp}\, \partial X ,
\]
i.e.\ supposons
\begin{itemize}
  \item[(i)] $Y \parallel \partial X$ \quad i.e.\ $\forall\, Y' < X$, $Y' \parallel Y$
  \item[(ii)] $Y \in \mathrm{supp}_{\mathcal{M}} X$ \quad i.e.\ $\forall\, Z \in \mathcal{M}$, $Z \parallel X \Rightarrow Z \parallel Y$.
\end{itemize}
Alors $Y \ll X$ \struck{\ill{}}.

\struck{\emph{Corollaire} Soient $F$ une figure, $G$ une sous-figure. Alors pour \struck{figures} $X \in \mathcal{M}$, les conditions suivantes sont équivalentes (i) $H \ll F$,}
\note{début de corollaire barré de longs traits obliques}

\emph{Corollaire}
\[
  X^{\circ} = \lbrace Y \in \mathrm{Omb}(X) \mid Y_{\partial X} = \emptyset ,\ \text{i.e.}\ \widetilde{Y} \cap \mathrm{Omb}(\partial X) = \emptyset \rbrace
\]
En effet, si $Y \ll X$, \struck{la relation \ill{} \ill{} $Y_{\partial X} = \emptyset$}, on sait par At D que
\[
  Y_{\partial X} = \emptyset \Longleftrightarrow Y \parallel \partial X
\]
(on a $\Longleftarrow$ sans hyp.\ At D) donc l'ens.\ du second membre \struck{\ill{}} est contenu dans $X^{\circ}$ par At supp 1. L'inclusion inverse, compte tenu de \uncertain{($\delta$)}, est triviale (NB $Y \in \mathrm{Omb}(X) \Longrightarrow Y \in \mathrm{supp}_{\mathcal{M}} X$).

\page{87}
\emph{Remarque} On peut renforcer \uncertain{un peu} At supp 1 ainsi :

\emph{At supp 1 (spéc)} Soient $X, Y$, \struck{$L'$} $\in \mathcal{M}$, \struck{\ill{}} $L \leq Y$. Supposons
\begin{itemize}
  \item[(i)] $\mathrm{supp}\, Y \subset X$
  \item[(ii)] $L \ll \partial X$
  \item[(iii)] « $Y \smallsetminus L$ » $\parallel \partial X$, de façon précise, $\forall\, Z \in \mathrm{Omb}(Y)$,
  \[
    Z|L = \emptyset \Longrightarrow Z \parallel \partial X
  \]
\end{itemize}
Alors $Y \ll X$ (et $L = Y_{\partial X}$)
\note{le $L$ est partout écrit en surcharge sur un autre signe (un $F$ ?, au-dessus duquel est écrit un $\widetilde{F}$ biffé) ; « $\mathrm{supp}\, Y \subset X$ » est sur la page, où l'on attend $\mathrm{supp}\, X$ ; dans (iii), un $\forall$ biffé précède « $Z \in \mathrm{Omb}(Y)$ »}

Ceci suggère qu'il doit exister une définition \uncertain{raisonnable} de la relation $\ll$, dans $\mathcal{M}$, au moyen de $\leq$ et de la relation $\parallel_{\mathcal{M}}$ — mais cette dernière relation \ill{} seulement pour deux \uncertain{doublets} de $\mathcal{M}$, $X, Y$, mais aussi pour des affirmations virtuelles
\[
  Y \smallsetminus L \parallel Z
\]
$Y, Z \in \mathcal{M}$, \struck{$L \in \widetilde{F}$, $L \leq Y$} $L$ partie fermée de $\widetilde{Y}$
\note{entre $\parallel$ et $Z$, un signe biffé}
voire même
\[
  Y \smallsetminus L \parallel Z \smallsetminus K
\]
— cette relation intervient \uncertain{implicitement} dans At C (spéc), mais circonscrite à l'\uncertain{intérieur} de \struck{$Z$} $S \ll Y \smallsetminus L$, $T \ll Z \smallsetminus K$ tels que\ldots

\page{88}
Faisons abstraction pour le moment des axiomes de supports, et de la définition proposée de $X^{\circ}$. Posons

\struck{\emph{Définition} (i) Soit $X \in \mathcal{M}$. Posons}
\struck{$\mathrm{Omb}^{\circ\circ}(X) = \lbrace Y \in \mathrm{Omb}(X) \mid Y \parallel \partial X \rbrace$}
\struck{(évidemment $\mathrm{Omb}^{\circ\circ}(X) \subset \mathrm{Omb}^{\circ}(X) = \mathrm{Omb}(X) \smallsetminus \mathrm{Omb}(\partial X)$)}

\struck{(ii) Posons}
\struck{$X^{\circ} \parallel Y^{\circ} \overset{\mathrm{def}}{\Longleftrightarrow} \mathrm{Omb}^{\circ\circ}(X) \cap \mathrm{Omb}^{\circ\circ}(Y) = \emptyset$}
\struck{i.e.\ il n'existe pas de $Z \in \mathcal{M}$} \struck{avec $Z \ll X, Y$, $Z|\partial X = \emptyset$, $Z|\partial Y = \emptyset$,} \struck{ou encore,} \struck{$Z \ll X, Y \Rightarrow{}$} \struck{$\exists\, Z' \in \widetilde{Z}$} \struck{avec $Z' < X$ ou $Z' < Y$.}
\note{la définition est encadrée et barrée de deux longs traits obliques ; devant « $\mathrm{Omb}^{\circ\circ}(X)$ », un « $X^{\circ}$ » surchargé ; dans la dernière formule, un signe biffé suit « $Z' <$ » ; on attend $Z' \leq \partial X$ ou $Z' \leq \partial Y$}

On aimerait avoir, pour $X, Y \in \mathcal{M}$
\[
  X \parallel Y \Longleftrightarrow \forall\, X' \in \widetilde{X},\ Y' \in \widetilde{Y},\ X'^{\circ} \parallel Y'^{\circ}
\]
On a toujours $\Longrightarrow$, car $X \parallel Y$ implique a fortiori $\nexists\, Z$ avec $Z \ll X, Y$ donc a fortiori $\nexists\, Z$ avec $Z \ll X'^{\circ}, Y'^{\circ}$ (si $X' \in \widetilde{X}$, $Y' \in \widetilde{Y}$). L'implication $\Longleftarrow$. Elle est conséquence plus délicate \ill{} de At L 2 et At L 3, \struck{\ill{}} \ill{}
\[
  X \parallel Y \Longleftrightarrow \forall\, x \in \mathrm{omb}(X),\ \forall\, y \in \mathrm{omb}(Y),\quad x \parallel y
\]
\note{le trait oblique qui barre la définition descend jusqu'au bas de la page, mais semble ne viser que celle-ci ; le reste est transcrit comme non biffé}

\page{89}
\emph{Définition} Soit $X \in \mathcal{M}$. Posons
\[
  X^{\circ} = \Big\lbrace Y \in \mathrm{Omb}(X) \ \Big|\ Y \parallel \partial X \ \text{i.e.}\ Y_{\partial X} = \emptyset \ \text{i.e.}\ \widetilde{Y} \cap \mathrm{Omb}(\partial X) = \emptyset \ \text{i.e.}\ \nexists\, Y' \leq Y,\ Y' \ll X' < X \Big\rbrace
\]
\note{le $\nexists$ est lu ainsi sans certitude (un $\exists$ n'est pas exclu) ; les trois gloses « i.e. » sont empilées à droite et rassemblées par une accolade ; devant la troisième, un mot biffé. Sous l'énoncé, à gauche, « \uncertain{ou} $\mathrm{Omb}^{\circ\circ}(X) \subset \mathrm{Omb}^{\circ}(X)$ », reste de la page 88}

Donc pour $X, Y \in \mathcal{M}$, la relation
\[
  X^{\circ} \parallel Y^{\circ}
\]
\struck{\ill{}} \struck{définie}, elle signifie
\[
  \forall\, X' \in X^{\circ},\ Y' \in Y^{\circ},\ \text{on a}\ X' \parallel Y' .
\]
\marginal{Attention, pour \struck{\ill{} \ill{}} que $X^{\circ} \cap Y^{\circ} = \emptyset \Longrightarrow X^{\circ} \parallel Y^{\circ}$, il faudrait pour $x, y \in \mathcal{L}$, $x \neq y \Longrightarrow x \parallel y$ ! (\ill{})}
\note{note encadrée d'un trait vertical, à droite de la formule}

On \uncertain{voudrait} prouver ceci (OK sous réserve \uncertain{de} At L2, L3)

\emph{Lemme 1} (?) (Soient $X, Y \in \mathcal{M}$. Alors
\[
  X \parallel Y \Longleftrightarrow \forall\, X' \in \widetilde{X},\ Y' \in \widetilde{Y},\ \text{on a}\ X'^{\circ} \parallel Y'^{\circ} .
\]
\note{le « 1 » est suivi d'un point d'interrogation entre parenthèses ; la parenthèse ouvrante qui précède « Soient » n'est pas refermée}

\emph{Dém.} $\Longrightarrow$ est clair, car $X'^{\circ} \subset \mathrm{Omb}(X') \subset \mathrm{Omb}(X)$, et de même $Y'^{\circ} \subset \mathrm{Omb}(Y') \subset \mathrm{Omb}(Y)$, et on a $\Longrightarrow$ \uncertain{par} At D$_0$. L'implication $\Longleftarrow$ est moins claire, \struck{\ill{} \ill{} \ill{} \ill{}} \add{elle \uncertain{utilise}} L2, L3.
\note{« Elle l'est » (?) souligné en marge droite, à la hauteur de « moins claire »}

\emph{Lemme 2} (At \uncertain{L1}, L2, L3) \struck{Supposons} Soit $Y \ll X$. Alors
\begin{itemize}
  \item[a)] $Y \overset{\circ}{\ll} X$ ssi $Y^{\circ} \subset X^{\circ}$. On a
  \item[b)] $Y \overset{\circ\circ}{\ll} X$ (i.e.\ $Y \in X^{\circ}$) ssi $\forall\, Y' \in \widetilde{Y}$, $Y'^{\circ} \subset X^{\circ}$.
\end{itemize}
\note{l'énoncé est marqué d'un trait vertical en marge ; « L1 » est entouré ; « ssi » est écrit « sss » ; « i.e. $Y \in X^{\circ}$ » est au-dessus de la ligne}
\marginal{\emph{NB} $Y \in X^{\circ}$ signifie aussi que $Y \ll X$, et l'application $\varphi : \widetilde{Y} \to \widetilde{X}$ est constante de valeur le plus grand élément de $\widetilde{X}$. \ill{} de $Y \in X^{\circ}$, \ill{} \ill{} \ill{} $Y \overset{\circ\circ}{\ll} X$ (raffinement \ill{} \ill{}) \ill{} $Y \in \mathrm{Omb}^{\circ\circ}(X)$, \ill{} $Y \overset{\circ}{\ll} X$ \ill{} \ill{} raffinement \emph{intérieur} de $X$, \ill{} \ill{} \ill{} $Y \in \mathrm{Omb}^{\circ}(X)$)}
\note{longue note oblique dans la marge gauche, en face des lemmes 1 et 2, lue en partie ; plus haut, une autre, « partie fermée de $\mathcal{M}$, $\ll$ », rattachée à la définition}

\emph{Dém} a) Supposons $Y \overset{\circ}{\ll} X$, prouvons $Y^{\circ} \subset X^{\circ}$. Soit $Y' \in Y^{\circ}$ i.e.\ $Y' \ll Y$, $Y' \parallel \partial Y$, donc $Y' \ll X$, il faut prouver $Y' \parallel \partial X$ ou encore $Y'_{\partial X} = \emptyset$.

\page{90}
\begin{tikzcd}
Y' \arrow[r, no head, "\ll" description] & Y \arrow[r, no head, "\overset{\circ}{\ll}" description] \arrow[d, no head, "\geq" description] & X \arrow[d, no head, "\geq" description] \\
 & Y_{\partial X} \arrow[r, no head, "\ll" description] \arrow[d, no head, "\geq" description] & \partial X \\
 & \partial Y &
\end{tikzcd}

\note{petit schéma en haut à gauche de la page : les signes verticaux sont des $\leq$ tournés, le terme le plus bas étant le plus petit}

Or $Y'|\partial X$ (induction dans $X$) $=$ \struck{\ill{}} $Y'|(Y_{\partial X})$ (induction dans $Y$) \struck{\ill{}} \struck{d'hypothèse}. Comme $Y_{\partial X} \leq \partial Y$ (par hyp.\ $Y \overset{\circ}{\ll} X$), on a
\[
  Y'|Y_{\partial X} \leq Y'|\partial Y
\]
\struck{\ill{}} or $Y'|\partial Y$ ou $Y'_{\partial Y} = \emptyset$
par hyp.\ $Y' \in Y^{\circ}$, d'où $Y'|Y_{\partial X} = \emptyset$ cqfd.

Supposons $Y^{\circ} \subset X^{\circ}$ prouvons $Y \overset{\circ}{\ll} X$, i.e.\ $Y \not\ll \partial X$. Je ferai usage ici
\[
  \boxed{Y^{\circ} \neq \emptyset}
\]
(ce qui résulte de \add{l'axiome} de divisibilité L1, \ill{} \ill{} \ill{} \ill{} : L1\ldots) Soit $Z \in Y^{\circ}$, donc $Z \in X^{\circ}$ par $Y^{\circ} \subset X^{\circ}$, donc $Z \parallel \partial X$. Si on avait $Y \ll \partial X$, on aurait $Z \ll \partial X$, d'où $Z \parallel Z$ par At D, absurde.
\note{le signe de « $Y \not\ll \partial X$ » est un $\ll$ traversé d'un trait, lu comme une négation, douteux ; l'argument qui suit en fait bien une négation}
\marginal{\emph{NB} Mais on n'a \ill{} \ill{} \ill{} hyp.\ \ill{} « $Y^{\circ} \neq \emptyset$ » \ill{} conditions \emph{équivalentes} : L1}
\note{note oblique en marge gauche, en face de l'encadré, lue en partie ; « équivalentes » est souligné}

b) Supposons $Y \in X^{\circ}$, prouvons $\forall\, Y' \in \widetilde{Y}$, $Y'^{\circ} \subset X^{\circ}$. En \struck{\ill{}} fait, on sait que $Y'^{\circ} \subset \mathrm{Omb}(Y')$ \ill{} $\subset X^{\circ}$, mais en fait $\mathrm{Omb}(Y') \subset \mathrm{Omb}(Y)$ \ill{}, mais en fait \struck{\ill{}} On a aussi
\[
  Y \in X^{\circ} \Longrightarrow \mathrm{Omb}(Y) \subset X^{\circ}
\]
i.e.\ $X^{\circ}$ est une partie de $\mathcal{M}$ \struck{\ill{}} fermée pour $\ll$, comme on s'en \uncertain{apercevait} déjà : Donc $\Longrightarrow$ dans b) est trivial. Prouvons $\Longleftarrow$, i.e.

\page{91}
supposons
\[
  \bigcup_{Y' \in \widetilde{Y}} Y'^{\circ} \subset X^{\circ} , \quad \text{prouvons}\ Y \in X^{\circ} ,
\]
i.e.\ $Y_{\partial X} = \emptyset$, \struck{Soit $Z$} i.e.\ $\nexists\, Z \in \widetilde{Y}$, avec $Z \ll \partial X$. En effet, pour un tel $Z$, \struck{on} \struck{la relation} On aurait par hyp.\ $Z^{\circ} \subset X^{\circ}$, donc $Z \overset{\circ}{\ll} X$ par a), d'où $Z \not\ll \partial X$. On gagne.

\emph{NB} Soient $X, Y \in \mathcal{M}$, \add{$X \neq Y$}. Alors
\[
  X^{\circ} \cap Y^{\circ} = \emptyset
\]
\note{« $X \neq Y$ » est écrit au-dessus d'une condition surchargée, peut-être $X \between Y$, et rattaché par un trait}
\marginal{Mieux, $X^{\circ} \parallel Y^{\circ}$ cf p.~95}
En effet, soit $L = X \cap Y$. Si on avait $Z \in X^{\circ} \cap Y^{\circ} \subset \mathrm{Omb}(X) \cap \mathrm{Omb}(Y) = \mathrm{Omb}(L)$, $Z \ll L$. Or on n'aurait \struck{donc} pas $L = X$, $L = Y$ (puisque $X \neq Y$), disons $L \neq X$ donc $L \subset \partial X$. \ill{} $Z \parallel \partial X$, $Z \parallel Z$, Donc $Z \ll \partial X$, \ill{} \ill{} absurde !

\struck{On a envie de} Soit maintenant $\Phi \subset \mathcal{M}$ \struck{\ill{}} une partie finie de $\mathcal{M}$, fermée \uncertain{pour} $\leq$, formée de \struck{figures} \add{multistrates} deux à deux compatibles. OPS $\Phi = \widetilde{F}$ (quitte à \struck{\ill{}} \ill{} $\mathfrak{F}$, \ill{} \ill{} partie de $\mathcal{M}$ — mais les

\page{92}
\uncertain{disposons} \ill{} actuelle \ill{} \uncertain{dispensée} de $I$ que par l'intermédiaire des strates $\mathcal{M}$, $\leq$, $\ll$, $\between$. \struck{Il sera commode de \ill{} \ill{} \ill{}, et \ill{} \ill{} \ill{} \ill{} \ill{} ordonné $I$, \ill{} \ill{} $\widetilde{F}$.}
\note{trois lignes barrées, lues par fragments}

Pour toute partie finie $A$ de $\widetilde{F} = \Phi$, posons
\[
  S_A = \mathrm{Supp}_{\mathcal{M}} \Big( \bigcup_{X \in A} X^{\circ} \Big) \in \Sigma_{\mathcal{M}}
  = \mathop{\mathrm{Sup}_{\Sigma_{\mathcal{M}}}}_{X \in A} \mathrm{Supp}(X^{\circ}) .
\]

(\emph{NB} Si on admet At supp.\ 1, on a $\mathrm{Supp}\, X^{\circ} = X^{\circ}$ i.e.\ $X^{\circ}$ est déjà une partie-support de $\mathcal{M}$\ldots)

Par définition et par associativité des Sup, si $A = \bigcup_i A_i$ on aura aussi
\[
  S_A = \mathop{\mathrm{Sup}_{\Sigma_{\mathcal{M}}}}_{i} S_{A_i} .
\]
On a ainsi une application
\[
  \mathfrak{P}(\widetilde{F}) \longrightarrow \Sigma_{\mathcal{M}} , \qquad A \longmapsto S_A
\]
de l'ens.\ ordonné des parties de $\widetilde{F}$ dans l'ens.\ \struck{\ill{}} ordonné des supports pour $\mathcal{M}$, application évidemment croissante.

\page{93}
J'aimerais prouver le

\emph{Théorème} (a) L'application précédente est injective, plus précisément elle induit un \emph{iso.\ d'ens.\ ordonnés} de $\mathfrak{P}(\widetilde{F})$ avec \struck{\ill{}} un sous-ensemble ordonné de $\Sigma_{\mathcal{M}}$. (b) Cette application est commutative aux Sup (trivial) et aux Inf quelconques. (c) On a, pour deux $A, B \in \mathfrak{P}(\widetilde{F})$, avec $B \subset A$,
\[
  S_{A \smallsetminus B} = S_A \cap \complement_{\mathcal{M}}(S_B)
\]
où $\complement_{\mathcal{M}}$ est l'anti-involution canonique de l'ens.\ des supports (\uncertain{complémentaire} \uncertain{support}). (d) \struck{On a} Si $A$ est une partie fermée de $\widetilde{F}$ (donc de la forme $\widetilde{G}$, où $G$ est une sous-figure de $F$) on a
\[
  S_A = \mathrm{supp}\, G \overset{\mathrm{def}}{=} \mathrm{Supp}(\hat{A}) .
\]
\note{les lettres (a) à (d) sont cerclées ; au début de la formule de (c), un mot biffé ; « iso.\ d'ens.\ ordonnés » est souligné ; « (trivial) » au-dessus de « Sup » ; sous (c), deux formules raturées d'une masse de traits, illisibles ; dans (d) la page porte « $\widetilde{C}$ », pour $\widetilde{G}$, et un petit chapeau (?) sur $A$, avec $\mathcal{M}$ en dessous}

\emph{NB} En fait, $S_A$ est défini dès qu'on se donne \struck{\ill{}} $F$ i.e.\ $\widetilde{F} = \Phi \subset \mathcal{M}$, i.e.\ la figure \ill{} à \ill{} $A \subset \mathcal{M}$ formé de multistrates mutuellement compatibles. Donc dans ce formulaire \ill{} \ill{}

\page{94}
formulaire ainsi, on travaille \uncertain{avec} \ill{} d'opérations, \ill{} sur des parties $A, B, A_i$ de $\mathcal{M}$, pas \uncertain{nécess.}\ fermées, et formées d'objets mutuellement compatibles. Commençons par d) (OPS $A = \widetilde{F}$)
\[
  S_{\widetilde{F}} \overset{\mathrm{def}}{=} \mathrm{Supp}\Big( \bigcup_{X \in \widetilde{F}} X^{\circ} \Big) = \mathrm{Supp}\, F
\]
ce qui revient, par définition des supports, à ceci, pour $\forall\, Y \in \mathcal{M}$
\[
  Y \parallel F \Longleftrightarrow \forall\, X \in \widetilde{F},\ Y \parallel X^{\circ}
\]
\[
  \big( Y \parallel \widetilde{F} \big)
\]
et cela résulte aussitôt du cor.\ au lemme 1.
\note{un « $\forall\, X$ » biffé précède le quantificateur ; « $(Y \parallel \widetilde{F})$ » est écrit sous « $Y \parallel F$ », relié par une double barre verticale ; le $F$ de « $\mathrm{Supp}\, F$ » et de « $Y \parallel F$ » est écrit en surcharge}

Voyons a) i.e.
\[
  S_A \subset S_B \Longrightarrow A \subset B
\]
Soit donc $X \in A$, on doit prouver $X \in B$, on est ramené à prouver que
\[
  S_{\lbrace X \rbrace} \subset S_B \Longrightarrow X \in B
\]
\[
  \big( S_{\lbrace X \rbrace} = \mathrm{supp}(X^{\circ}) \big)
\]
i.e.

\emph{Lemme 3} Soient $X \in \mathcal{M}$, $B \subset \mathcal{M}$, tels que $B \cup \lbrace X \rbrace$ soit admissible (formé d'objets mutuellement \struck{admissibles} \uncertain{compatibles}). Alors
\[
  X^{\circ} \subset S_B \Longrightarrow X \in B
\]
\note{le « 3 » est repassé ; « soit admissible » est surmonté d'un mot court}
\marginal{\emph{NB} Cela \struck{\ill{}} \ill{} \ill{} \ill{} $B$ \ill{} \ill{} $S_B$ \ill{}}
\note{note oblique en marge gauche, lue en partie}

La preuve \struck{\ill{}, \ill{} \ill{}} consiste \ill{} $\forall\, Z$ \struck{\ill{}} tel que $Z \parallel \bigcup_{Y \in B} Y^{\circ}$, on a $Z \parallel X^{\circ}$

\page{95}
donc tout revient à montrer ceci :
\[
  X \notin B \Longrightarrow \exists\, Z \in \mathcal{M},\ Z \parallel \bigcup_{Y \in B} Y^{\circ} \ \text{et}\ Z \overline{\parallel}\, X^{\circ}
\]
et pour ceci, il suffit de prendre
\[
  Z \in X^{\circ}
\]
(\struck{\ill{}} existe par L1, impliquant $X^{\circ} \neq \emptyset$)
et appliquer le
\note{la barre est tracée au-dessus de $\parallel$ ; le mot biffé est sans doute « existe »}

\emph{Lemme 4} Soient $X, Y \in \mathcal{M}$, $X \between Y$, $X \neq Y$. Alors $X^{\circ} \parallel Y^{\circ}$ (a fortiori, $X^{\circ} \cap Y^{\circ} = \emptyset$)
\note{le lemme est marqué de trois traits verticaux dans la marge}

Soient $X' \in X^{\circ}$, $Y' \in Y^{\circ}$, il faut prouver $X' \parallel Y'$. Soit $L = X \cap Y$. Par At D, il suffit \uncertain{de prouver} $X'_L \parallel Y'_L$. Si $L \leq \partial X, \partial Y$, \struck{on} i.e.\ $L \neq X, Y$, on a
\[
  X'_L \leq X'_{\partial X} = \emptyset , \qquad Y'_L \leq Y'_{\partial X} = \emptyset
\]
\struck{et ainsi} on a donc $X'_L = Y'_L = \emptyset$, OK. Si \struck{\ill{}} \struck{supp}, $L =$ \struck{$X$} $Y$ i.e.\ $Y \leq X$. Comme $Y \neq X$, on a $Y \leq \partial X$ \ill{}, $Y' \leq Y \leq \partial X$, et comme $X' \parallel \partial X$ on a $X' \parallel Y'$, OK.
\note{« $Y'_{\partial X}$ » est sur la page, où l'on attend $Y'_{\partial Y}$ ; le premier $\leq$ de la ligne est très appuyé et pourrait être un $\ll$}

Pour prouver b), \struck{\ill{}} (\uncertain{pour le} \uncertain{sup}) on est ramené à prouver la \struck{relation} compatibilité \ill{} c) pour les complémentaires. \struck{Or}

\page{96}
En fait, le th.\ est un \add{cas particulier} \struck{conséquence} de l'énoncé formel que voici :
\marginal{? \quad Faux}
\note{un grand point d'interrogation en marge gauche en face de « Lemme », et plus bas, en biais, « Faux » (suivi d'une petite croix)}

\emph{Lemme} Soit $\mathcal{M}$ un ens.\ muni d'une relation sym.\ antiréflexive $\parallel$, d'où une notion de support, \uncertain{supports} pour une partie de $\mathcal{M}$, et l'ensemble $\Sigma_{\mathcal{M}} = \Sigma_{\mathcal{M}, \parallel}$\note{l'indice $\mathcal{M}$ du premier $\Sigma$ est biffé} des parties-supports, avec $\mathrm{Inf} = \bigcap$, Sup, et l'anti-involution $\complement$ canonique (le complément) — (grâce à la symétrie \add{de $\parallel$}) satisfait (plus subtil !)
\[
  S \cap \complement S = \emptyset \ \ (\text{plus petit élt.\ de}\ \Sigma)
\]
\note{après le signe $=$, un $\emptyset$ biffé, et le mot « plus petit » récrit au-dessus}
(grâce à l'antiréflexivité) et \struck{\ill{} \ill{} $\Sigma = \mathrm{cosupp}\, \mathcal{M} = \lbrace x \in \mathcal{M}$}
\[
  S \vee \complement S = \mathcal{M} \qquad (\text{plus grand élt de}\ \Sigma) .
\]
Soit $I$ un ens.\ d'indices, et $(A_i)_{i \in I}$ une \add{\uncertain{famille}} de parties de $\mathcal{M}$ (en l'occurrence, $I = \widetilde{F}$, $A_i = X_i^{\circ}$). On suppose
\begin{itemize}
  \item[a)] $\forall\, i$, $A_i \neq \emptyset$.
  \item[b)] $\forall\, i, j \in I$, $i \neq j$, $A_i \parallel A_j$.
\end{itemize}
Pour toute partie $J$ de $I$, soit
\[
  S_J = \mathop{\mathrm{Sup}_{\Sigma}}_{i \in J} \mathrm{supp}\, A_i
\]
Alors l'application
\[
  \varphi : J \longmapsto S_J : \mathfrak{P}(I) \longrightarrow \Sigma
\]
\struck{est injective} induit un \uncertain{iso.}\ de $\mathfrak{P}(I)$ sur un sous-ens.\ ordonné \add{$\Sigma'$} de $\Sigma$, elle commute aux Sup et aux Inf quelconques (donc $\Sigma'$ est stable dans $\Sigma$ pour ces opérations), enfin pour $A, B \in \mathfrak{P}(I)$, $B \subset A$, on a
\[
  S_{A \smallsetminus B} = S_A \cap \complement S_B .
\]
\note{« est injective », biffé, précède « induit » ; « $\Sigma'$ » est écrit au-dessus de la ligne ; « (en l'occurrence, \ldots) » est écrit sous la ligne ; dans la dernière formule, les indices $A$, $B$ sont écrits en surcharge}

\page{97}
\emph{Dém} Commençons \uncertain{par l'injectivité}.

Soient $A, B \subset I$, on prouve
\[
  S_A \subset S_B \Longrightarrow A \subset B
\]
i.e.\ $i \in A \Rightarrow i \in B$. Cela revient au \struck{lemme}

\emph{Lemme 1} Si $i \in I$, \struck{\ill{} \ill{}} $B \subset I$, alors
\[
  S_{\lbrace i \rbrace} \subset S_B \Longrightarrow i \in B
\]
Résulte du

\emph{Lemme 2} Si \struck{$A$}, $B$, $C \subset I$ avec $B \cap C = \emptyset$, alors
\[
  S_B \parallel S_C \qquad \text{a fortiori}\ S_B \cap S_C = \emptyset
\]
Immédiat par condition b). Le lemme 1 en résulte — car si $i \notin B$ on aurait $S_i \parallel S_B$, donc (comme $S_i \subset S_B$) $S_i \parallel S_i$ donc $S_i = \emptyset$, a fortiori $A_i = \emptyset$, contrairement à a).
\note{« $B \cap C = 0$ » sur la page, pour $\emptyset$}

La commutation de $\varphi$ aux Sup est triviale sur la définition, pour la commutation aux Inf, par passage aux complémentaires, on est ramené à vérifier \uncertain{p.r.\ à} $S_I$, \ill{} on est ramené \ill{} la formule sur $S_{A \smallsetminus B}$. Mais posons
\[
  A = B \cup B' \qquad \text{où}\ B' = A \smallsetminus B
\]
Donc
\[
  S_A = S_B \vee S_{B'}
\]
et la formule s'écrit :
\[
  S_{B'} = (S_B \vee S_{B'}) \cap \complement S_B ,
\]
i.e.\ on est ramené au
\note{les lettres $A$, $B$ de ce passage sont repassées en gras}

\emph{Lemme 3} Soient \struck{$A, B \in$} $S, S' \in \Sigma$, avec $S \parallel S'$, alors
\[
  S' = (S \vee S') \cap \complement S \qquad \big( = \lbrace X \in S \vee S' \mid X \parallel S \rbrace \big)
\]
\marginal{Faux}

\emph{Dém} L'inclusion $S' \subset (S \vee S') \cap \complement S$ est triviale, la relation $S \parallel S'$ étant équivalente à : $S' \subset \complement S$. En sens inverse, si $S \vee S' = \mathcal{M}$, ça reviendrait à prouver que $S$, $S'$ sont complémentaires l'un de l'autre.

\page{98}
Mais c'est faux ! Prenons p.\ ex.\ un graphe \ill{} ayant une partie
\drawing{dessin : une chaîne de quatre sommets reliés $A'$ — $A$ — $B$ — $B'$, des pointillés prolongeant le graphe au-delà de $A'$ et de $B'$}
($A$ et $B$ ne sont liés qu'aux sommets $A'$, $B'$ et entre eux, comme on l'indique) On trouve
\[
  \mathrm{Cosupp}\, A = \lbrace A', B \rbrace , \qquad \mathrm{supp}\, A = \mathrm{cosupp}\, \lbrace A', B \rbrace = A'
\]
et de même
\[
  \mathrm{Cosupp}\, B = \lbrace A, B' \rbrace , \qquad \mathrm{supp}\, B = B
\]
\note{« $\mathrm{supp}\, A = \mathrm{cosupp} \lbrace A', B \rbrace = A'$ » et « $\mathrm{supp}\, B = B$ » sont ainsi sur la page ; on attend sans doute $\mathrm{supp}\, A = A$ ; le « $A'$ » final est écrit en surcharge}
Comme $A \parallel B$, considérons \uncertain{les} $A$ et $B$ fermés,
\[
  \text{Comme}\quad A \cup B = \complement A \cap \complement B = \emptyset \qquad \text{pas de sommet lié à $A$ et $B$ à la fois !}
\]
donc
\[
  \mathrm{Supp}\, A \cup B = A \vee B = \mathcal{M} .
\]
\note{« $A \cup B = \complement A \cap \complement B$ » : on attend $\complement(A \cup B)$ ; « pas de sommet \ldots\ à la fois ! » est écrit à droite, sur deux lignes}
Le lemme affirme donc que $A = \complement B$, $B = \complement A$, ce qui est faux ! Rideau\ldots

Il semblerait donc que \add{la \uncertain{sollicitation} de l'intuition} \ill{} \ill{} \uncertain{figures} d'hyper\ill{} \uncertain{hypothétiques} \uncertain{tentaient} \ill{} \uncertain{dégager} \ill{} \ill{} délicates, qui \ill{} \uncertain{d'essence} telles que la formule pertinente
\[
  S_{A \smallsetminus B} = S_A \cap \complement S_B
\]
\ill{} qui serait un peu brutal ! Pourtant, je suppose qu'elle devrait être vraie \struck{en tous cas} \add{\uncertain{du moins}} sous l'hypothèse de divisibilité At L1 (avec l'approche que j'ai
\note{fin de page écrite vite, lue par fragments ; plusieurs insertions interlinéaires}

\page{99}
prise pour donner un sens intrinsèque aux « supports constructibles », via les $X^{\circ}$.

\subsection{26 juin}
\marginal{26 juin}
\note{date de sa main, soulignée, en tête du second alinéa de la page 99}

Je \uncertain{vois} peut-être plus simple, \struck{puisque} pratiquement je suis obligé ici \struck{de} (pour une façon de définir $X^{\circ}$ pour « support » au sens $\parallel$) \struck{de fait} d'admettre At L1 (en plus de At L2, At L3), autant la définir plus simplement comme
\[
  X^{\circ} = \mathrm{Supp}\, \mathrm{omb}(X)^{\circ}
\]
c'est pareil grâce au cor.\ p.~86 et au fait \struck{\ill{} $Y \in \mathrm{Omb}$ \ill{} \ill{}} \ill{} At L2, que pour \struck{\ill{}} partie-support $S$ de $\mathcal{M}$, on a
\[
  S = \mathrm{supp}(S \cap \mathcal{L}) . \qquad (\text{plus gén., si $S$ est fermée pour $\ll$, $\mathrm{supp}\, S = \mathrm{Supp}(S \cap \mathcal{L})$}) .
\]

Notons (pour nous rassurer)

\emph{Lemme} Soit \struck{$F$} $F \preccurlyeq X$, \struck{\ill{}} et considérons
\[
  \varphi : \widetilde{F} \longrightarrow \widetilde{X} , \quad \text{soit} \quad \widetilde{F}^{\circ} = \varphi^{-1}(\lbrace X \rbrace) = \lbrace X' \in \widetilde{F} \mid X' \overset{\circ}{\ll} X \rbrace .
\]
Alors
\[
  X^{\circ} = \mathop{\mathrm{Sup}_{\Sigma}}_{X' \in \widetilde{F}^{\circ}} X'^{\circ}
  \quad \Big( = \mathrm{Supp} \bigcup_{X' \in \widetilde{F}^{\circ}} X'^{\circ} = \mathrm{Supp} \bigcup_{X' \in \widetilde{F}^{\circ}} \mathrm{omb}(X')^{\circ} \Big)
\]
\note{le lemme est marqué de deux traits verticaux ; dans la parenthèse, un « $= \mathrm{Supp} \bigcup$ » biffé avant la dernière égalité ; « i.e. $X'^{\circ} \subset X^{\circ}$ » est écrit au-dessus de « $X' \overset{\circ}{\ll} X$ »}

Pourtant, dans le cas \uncertain{quasi-ensembliste}, on n'a \uncertain{pas} en général \ill{} $\mathrm{omb}(X)^{\circ} = \bigcup_{X' \in \widetilde{F}^{\circ}} \mathrm{omb}(X')^{\circ}$.

\struck{Il} L'inclusion $\supset$ est évidente du fait que les $X'^{\circ}$ sont $\subset X^{\circ}$. Prouvons $\subset$. \struck{Donc} Ceci équivaut au

\page{100}
\emph{Corollaire} $\forall\, Z \in \mathcal{M}$ tel que \add{$\forall\, X' \in \widetilde{F}^{\circ}$, on ait} $Z \parallel{}$ \struck{$X$} $\mathrm{omb}(X')^{\circ}$ \add{\uncertain{soit}}. \struck{Alors résulte} \uncertain{Alors} $Z \parallel \mathrm{omb}(X)^{\circ}$.
\note{le corollaire est marqué de deux traits verticaux, et d'un grand point d'interrogation en marge gauche ; « $\forall\, X' \in \widetilde{F}^{\circ}$, on ait » est écrit au-dessus de la ligne}

\struck{On} \add{Par At L2} On est ramené au cas où $Z = x \in \mathcal{L}$.

Je rappelle que $F \preccurlyeq X$ signifie
\[
  F \ll X \quad \text{et} \quad \mathrm{supp}\, F = \mathrm{supp}\, X .
\]
Si on a une fonction \add{$X \mapsto |X|$} \struck{support} à valeurs dans un ens.\ $P$ (cf p.~80, condition Supp$_P$ 1 — on suppose Supp$_P$ 2) \uncertain{alors} pour $F \ll X$, \ill{} \ill{} \ill{} la condition $F \preccurlyeq X$ signifie simplement que
\[
  |F| = |X| .
\]
Je dis que
\begin{itemize}
  \item[(a)] $|X^{\circ}| = |X| \smallsetminus |\partial X|$, et que
  \item[(b)] $\widetilde{F}^{\circ} = \lbrace X' \in \widetilde{F} \mid |X'^{\circ}| \subset |X^{\circ}| \rbrace = \lbrace X' \in \widetilde{F} \mid |X'^{\circ}| \cap |X^{\circ}| \neq \emptyset \rbrace$
\end{itemize}
\note{(a) et (b) sont cerclés ; devant (a) un « supp » biffé, devant (b) un « $\widetilde{F}$ » biffé ; les deux énoncés sont réunis par une accolade}
\marginal{On doit \ill{} \ill{} supposer \ill{}}
\note{note oblique en marge gauche, en face de « $F \ll X$ », lue en partie}

Comme $X^{\circ} \parallel \partial X$, \struck{\ill{}} et $X^{\circ} \subset{}$ \struck{$X$} $\mathrm{supp}\, X$, $|X^{\circ}| \subset |X| = |\mathrm{supp}\, X|$, et $|X^{\circ}| \cap |\partial X| = \emptyset$, donc
\[
  |X^{\circ}| \subset |X| \smallsetminus |\partial X|
\]
\note{la page s'arrête sur cette ligne, suivie de deux traits horizontaux}

\end{document}
